\chapter{Diseño de la arquitectura}
\label{chap:diseno}

\section{Criterios de diseño}

CitaLab separa la interpretación de una solicitud de su ejecución. El modelo
identifica lo que pide la persona y prepara la respuesta. El programa comprueba
los permisos, consulta la agenda y realiza las operaciones. Entender una petición
no concede autorización para ejecutarla.

Los requisitos del capítulo~\ref{chap:dominio-regulado} se aplican a las tres
arquitecturas. El acceso al perfil depende de la finalidad de la gestión. Las
operaciones con efectos necesitan una propuesta mostrada y una confirmación
expresa. El asistente trabaja con datos sintéticos y no realiza valoraciones
clínicas. Estas condiciones se mantienen al comparar M5, S1 y H.

Una interpretación incompleta puede resolverse preguntando por el dato que falta.
Si el modelo redacta una respuesta que pierde un dato importante, por ejemplo la
hora de la cita propuesta, CitaLab no muestra esa redacción. En su lugar muestra
un texto fijo que el programa compone con los datos leídos directamente de la
agenda, como \enquote{Propuesta de reserva: Pediatría, martes 3 de noviembre de
2026 a las 11:00 (presencial, Centro de Salud Demo). Todavía no se ha aplicado
ningún cambio. \textquestiondown{}Lo confirmas?}.
En cambio, una petición sin permisos o sin confirmación no puede producir una
escritura. La Tabla~\ref{tab:req-mecanismos} relaciona estas
decisiones con los requisitos del dominio. El Apéndice~\ref{app:requisitos-normativos}
recoge su trazabilidad normativa.

\begin{table}[!htbp]
\centering
\caption{Requisitos del dominio y mecanismo de la arquitectura que los satisface.}
\label{tab:req-mecanismos}
\small
\begin{tabularx}{\textwidth}{L{0.15\textwidth}Y}
\toprule
\textbf{Requisito} & \textbf{Mecanismo} \\
\midrule
REQ-CITA-01 & Identidad, paciente y representación tomados de la sesión. Regla \texttt{CMP-REPRESENT-01} en las consultas, en las propuestas sobre citas existentes y al confirmar. \\
REQ-CITA-02 & Servicios, huecos y citas solo a través de la capa de herramientas. \\
REQ-CITA-03 & Propuesta pública previa y confirmación explícita (\texttt{CMP-CONFIRM-01}). \\
REQ-CITA-04 & Revalidación de versión y disponibilidad bajo bloqueo. Idempotencia. \\
REQ-CITA-05 & Búsqueda restringida al servicio, fecha, hora, franja y modalidad pedidos. Sin sustitución silenciosa. \\
REQ-CITA-06 & Reprogramación atómica que comprueba el destino antes de liberar el origen. \\
REQ-CITA-07 & Vistas del perfil por finalidad. Las preferencias solo ordenan o adaptan la forma. \\
REQ-CITA-08 & Hechos de propuestas y recibos generados desde la agenda. Comprobación de la redacción. \\
REQ-CITA-09 & Avisos clínicos, rechazo fuera de alcance y parada ante emergencia (112). \\
REQ-CITA-10 & Modo de datos sintético obligatorio (\texttt{CMP-MODE-01}). \\
REQ-CITA-11 & Traza mínima \texttt{AuditEvent} separada del historial. Diagnóstico solo con activación expresa. \\
REQ-CITA-12 & Fuera del prototipo: procedimiento organizativo. \\
\bottomrule
\end{tabularx}
\end{table}

Las herramientas comprueban el acceso y las condiciones de escritura. Las
respuestas se revisan antes de mostrarlas para detectar pérdidas de información
o anuncios de efectos que no se han producido. Estas comprobaciones reducen
errores concretos, pero no bastan para valorar toda la conversación. Los fallos
de proveedor se registran aparte.

\section{Planos y fronteras de autoridad}

La arquitectura separa tres planos (Figura~\ref{fig:diseno-arquitectura}). El
plano conversacional interpreta y redacta. El plano de control concentra las
políticas, la confirmación y la ejecución. El plano de datos conserva la agenda
FHIR y los perfiles sintéticos. CitaLab, el portal web, da acceso a la
conversación y a la inspección de cada decisión.

\begin{figure}[!htbp]
\centering
\begin{tikzpicture}[
  node distance=7mm,
  role/.style={draw=ETSBlue,rounded corners,align=center,minimum height=9mm,
    text width=112mm,fill=ETSLightBlue!14,font=\small},
  control/.style={role,draw=ETSGold!75!black,fill=ETSGold!15},
  data/.style={role,draw=USRed,fill=USRed!7,text width=53mm},
  exchange/.style={{Latex[length=2mm]}-{Latex[length=2mm]},thick,draw=ETSBlue}
]
\node[role] (agents) {Plano conversacional: M5 (cinco agentes), S1 (agente único) y H (flujo programado)};
\node[control,below=of agents] (core) {Plano de control: capa de herramientas, motor de políticas,\\propuestas, confirmación, revalidación e idempotencia};
\node[data,below=12mm of core,xshift=-30mm] (portal) {Portal CitaLab\\conversación, inspección y evaluación};
\node[data,right=6mm of portal] (his) {Simulador HIS/FHIR\\agenda y perfiles sintéticos};
\draw[exchange] (agents) -- node[right,font=\scriptsize]{acciones tipadas / observaciones} (core);
\draw[exchange] (portal) -- (core);
\draw[exchange] (core) -- (his);
\end{tikzpicture}
\caption{Planos de la arquitectura. Los agentes solo actúan a través de la capa de herramientas.}
\label{fig:diseno-arquitectura}
\end{figure}

La frontera principal es la \textbf{capa de herramientas}: el único punto por
el que los modelos pueden pedir datos o cambios. Ofrece nueve operaciones
(Tabla~\ref{tab:herramientas}). Las que consultan citas, huecos o el perfil y
las que proponen cambios sobre una cita existente invocan el motor de
políticas, y \texttt{confirm} lo vuelve a invocar antes de escribir. El
programa también lee datos fuera de esas operaciones para preparar el contexto
de los agentes: un resumen de las citas del paciente de la sesión y el estilo
de respuesta guardado en su perfil, salvo que la persona haya renunciado a
usarlo. Esas lecturas se limitan al paciente que fija la sesión. Ningún modelo
recibe acceso a la base de datos, a consultas arbitrarias ni a código
ejecutable.

\begin{table}[!htbp]
\centering
\caption{Capa de herramientas compartida por M5, S1 y H.}
\label{tab:herramientas}
\small
\begin{tabularx}{\textwidth}{L{0.27\textwidth}Y}
\toprule
\textbf{Herramienta} & \textbf{Efecto y control} \\
\midrule
\texttt{list\_appointments} & Lista las citas propias de la persona autenticada, con filtros opcionales. Nunca devuelve citas ajenas. \\
\texttt{find\_slots} & Busca huecos libres de un servicio y los ordena (hora exacta, franja preferida autorizada, cronología). No reserva. \\
\texttt{get\_preferences} & Devuelve la vista del perfil para una finalidad (ordenar huecos, adaptar respuestas o canal de recordatorio). Se deniega si la persona renunció al perfil. \\
\texttt{propose\_booking}, \texttt{propose\_reschedule}, \texttt{propose\_cancel}, \texttt{propose\_reminder} & Crean una propuesta pública con todos sus datos, sin escribir nada. Una cita ajena se deniega sin revelar su contenido. \\
\texttt{confirm} & Única ruta de escritura: ejecuta la propuesta vigente si la persona la acepta con el botón Confirmar o con un mensaje afirmativo después de haberla visto. \\
\texttt{discard} & Descarta la propuesta pendiente. \\
\bottomrule
\end{tabularx}
\end{table}

\section{Confirmación y escritura}
\label{sec:diseno-confirmacion}

Se distinguen tres actos que suelen confundirse. La \emph{autenticación}
identifica la sesión. La \emph{autorización} decide si esa sesión puede leer o
escribir sobre un paciente. La \emph{confirmación} expresa que la persona acepta
una propuesta concreta. Toda gestión con efectos sigue dos tiempos: primero se
presenta una propuesta con servicio, fecha, hora, modalidad y centro, y después
la persona la confirma.

La confirmación se acepta por dos vías: el botón de la tarjeta de propuesta o
un mensaje afirmativo. En el segundo caso, la arquitectura debe interpretar el
mensaje como aceptación, y el código exige además que la propuesta se haya
mostrado en un turno anterior, que siga vigente y que el mensaje contenga una
expresión afirmativa sin ser una pregunta ni un condicional
(\enquote{\textquestiondown{}qué pasaría si confirmo?} nunca confirma). Si alguna
condición falla, el asistente vuelve a pedir la confirmación. Esta comprobación
del texto es léxica y sencilla. Rechaza los mensajes que empiezan por una
negación, como \enquote{espera, todavía no}, pero no una negación posterior,
como la de \enquote{de acuerdo, pero no lo confirmes}. Ante un mensaje así, la
protección depende de que la arquitectura no lo interprete como una aceptación.
El botón no tiene este límite, y en la campaña todas las escrituras se
confirmaron con él (apartado~\ref{sec:resultados-seguridad}).

Al ejecutar una reserva, un cambio o una cancelación, el simulador revalida la
versión del hueco y de la cita bajo un mismo bloqueo y aplica la operación de
forma atómica. En un recordatorio solo comprueba que la cita siga activa y sea
de la persona, así que el aviso se calcula sobre la hora que tenga la cita al
confirmar. Una clave de idempotencia
ligada a la propuesta impide duplicar efectos si la petición se repite. Si otra
persona ocupó el hueco entre la propuesta y la confirmación, la operación se
rechaza, la cita original se conserva y el asistente lo comunica sin reservar
otra alternativa por iniciativa propia.

\section{Identidad y representación}

La sesión fija tres datos que ningún mensaje puede cambiar: quién actúa, sobre
qué paciente y en nombre de quién. Una persona puede gestionar sus propias citas
y, si consta una relación de representación sintética, las del paciente
representado. Afirmar en la conversación que se actúa por otra persona no
concede ningún permiso: la capa de herramientas deniega el acceso con la regla
\texttt{CMP-REPRESENT-01} sin revelar el contenido del recurso ajeno.

\section{Perfil, memoria y personalización}

CitaLab no envía al modelo el historial completo para personalizar. Conserva
por separado el \emph{perfil} (preferencias entre conversaciones) y el
\emph{estado de la gestión} (lo que la persona pide ahora, las opciones
mostradas y la propuesta pendiente).

Cada atributo del perfil incluye valor, origen (declarado, importado o
inferido), confianza, vigencia, finalidad y procedencia. La resolución es
determinista: descarta lo retirado, caducado, no autorizado para la finalidad o
con confianza inferior a 0,5, y prioriza declarado sobre importado y este sobre
inferido. El valor 0,5 se fijó como un punto de corte sencillo: por debajo de él
una preferencia se considera demasiado incierta para usarla. Es una decisión de
diseño de este prototipo, no el resultado de un estudio, y puede cambiarse
modificando un único parámetro sin afectar al resto del sistema.

Si después de esos filtros quedan dos valores distintos con el mismo origen,
confianza y fecha (por ejemplo, mañana y tarde declaradas el mismo día), el
sistema no elige uno al azar. Ese atributo no se aplica: los huecos se muestran
por orden de fecha, como si no hubiera preferencia guardada, y la persona elige
o indica la franja que quiere. Los agentes reciben ese conflicto señalado en la
vista de preferencias.

Las preferencias se aplican con una precedencia explícita: \emph{petición actual
$>$ ajuste de la conversación $>$ perfil}. Una franja pedida en el mensaje
sustituye a la guardada. En M5 y H, pedir \enquote{no uses mis preferencias}
activa una marca que impide leer el perfil durante esa conversación. S1 no
tiene agente de contexto que la active, así que depende de que el modelo deje
de consultarlo. El control equivalente del portal activa la marca en las tres
arquitecturas. La Tabla~\ref{tab:preferencias-alcance}
resume para qué sirve cada preferencia y qué efecto puede tener.

Los ajustes que la persona hace durante una conversación no se escriben en su
perfil. Para conservarlos en futuras conversaciones, la persona abre el panel de
preferencias de la conversación en CitaLab y elige guardarlos en el perfil.
CitaLab le muestra entonces una propuesta con los valores nuevos, y solo cuando
pulsa Confirmar se guarda una nueva versión del perfil con esas preferencias
como declaradas. Si la rechaza, el perfil no cambia.

\begin{table}[!htbp]
\centering
\caption{Preferencias, finalidad y efecto.}
\label{tab:preferencias-alcance}
\small
\begin{tabularx}{\textwidth}{L{0.22\textwidth}L{0.22\textwidth}Y}
\toprule
\textbf{Atributo} & \textbf{Finalidad} & \textbf{Efecto} \\
\midrule
Franja horaria & Ordenar huecos & Ordena candidatos compatibles. Nunca elimina una restricción expresa ni crea disponibilidad. \\
Modalidad y centro & Ordenar huecos & No se aplican en esta versión, porque las agendas solo ofrecen citas presenciales en un centro. Una modalidad pedida que no existe no se sustituye por otra. \\
Idioma, detalle y tratamiento & Adaptar la respuesta & Cambian la forma, no los hechos ni la necesidad de confirmar. \\
Canal de recordatorio & Recordatorios & Canal por defecto de un recordatorio simulado si la persona no indica otro. \\
\bottomrule
\end{tabularx}
\end{table}

\section{Traza mínima e historial de investigación}

Cada consulta de citas, escritura, conflicto o denegación genera un registro
\texttt{AuditEvent} de FHIR con el actor, el paciente, la operación, la regla
aplicada, el resultado y la marca temporal. No incluye mensajes ni contenido de
la conversación y se conserva aparte del historial. Las búsquedas de huecos y
las lecturas del perfil quedan en las trazas de la conversación, no en esta
traza mínima. El historial de
investigación de CitaLab guarda las conversaciones sintéticas y las métricas de
cada llamada al modelo. Las entradas y salidas completas solo se conservan si
se activa expresamente el diagnóstico, con un plazo propio (7 días por defecto)
que no puede superar el del historial (30 días por defecto). En ningún caso se
guardan razonamientos internos del modelo.

\section{Respuestas: hechos del código, redacción del modelo}

Cada respuesta se construye a partir de un \emph{plan}: el tipo de mensaje
(opciones, propuesta, recibo, conflicto, listado, aclaración, información,
rechazo o escalado), los hechos que debe comunicar, que el código formatea a
partir de la agenda, los avisos obligatorios y, si procede, una pregunta. El modelo redacta ese plan con el idioma, el nivel de detalle y el
tratamiento aplicables.

Antes de mostrar la redacción se hacen cuatro comprobaciones deliberadamente
simples: que contenga las horas del plan, que no añada horas ajenas a la
petición, al plan o a la agenda de la persona, que mantenga la pregunta si la
hay y que no anuncie una reserva, cancelación o cambio cuando el plan no es un
recibo.
Si alguna falla, se muestra el texto determinista del plan. Los avisos
normativos (valoración clínica, emergencia, datos de terceros) se insertan por
código y no dependen de la redacción. El
apartado~\ref{sec:resultados-respaldos} analiza cuántas veces actuaron estas
comprobaciones en la campaña y con qué acierto.

\section{Las tres arquitecturas}

Las tres variantes comparten el núcleo descrito y difieren en cómo producen las
decisiones semánticas. Con la terminología del capítulo~\ref{chap:fundamentos},
H es un \emph{workflow}: el programa fija la ruta. M5 y S1 dejan al modelo
elegir el siguiente paso entre acciones acotadas mediante un ciclo de acción y
observación, sin generar ni conservar razonamientos intermedios. La
Tabla~\ref{tab:agentes-logicos} indica qué agentes lógicos tiene cada variante
y cómo se ejecutan, y la Tabla~\ref{tab:comparadores-mecanismos} compara quién
toma cada tipo de decisión.

\begin{table}[!htbp]
\centering
\caption{Agentes lógicos y llamadas por arquitectura.}
\label{tab:agentes-logicos}
\small
\begin{tabularx}{\textwidth}{L{0.10\textwidth}L{0.33\textwidth}Y}
\toprule
\textbf{Variante} & \textbf{Agentes lógicos (rol, instrucciones y contrato)} & \textbf{Nodos y procesos} \\
\midrule
M5 & Contexto, cumplimiento, perfiles, orquestación y personalización & Seis nodos LangGraph en un único proceso. El mismo modelo para los cinco roles \\
S1 & Agente único con nueve herramientas & Un bucle de llamadas a funciones en el mismo proceso \\
H & Interpretación y redacción (tareas LLM acotadas) & Secuencia programada en código \\
\bottomrule
\end{tabularx}
\end{table}

\begin{table}[!htbp]
\centering
\caption{Mecanismos de M5, S1 y H bajo controles comunes.}
\label{tab:comparadores-mecanismos}
\small
\begin{tabularx}{\textwidth}{L{0.17\textwidth}YYY}
\toprule
\textbf{Aspecto} & \textbf{M5} & \textbf{S1} & \textbf{H} \\
\midrule
Interpretación & Agente de contexto & El propio agente & Agente de contexto (tarea LLM) \\
Riesgos normativos & Agente de cumplimiento y reglas & Instrucciones y reglas & Solo reglas \\
Preferencias & Agente de perfiles decide cuáles aplicar & Herramienta \texttt{get\_preferences} & Reglas de precedencia \\
Siguiente paso & Orquestador LLM entre acciones habilitadas & Bucle de llamadas a herramientas & Secuencia programada \\
Redacción & Agente de personalización & El propio agente & Tarea LLM de redacción \\
\midrule
Controles y efectos & \multicolumn{3}{p{0.68\textwidth}}{Mismas herramientas, políticas, vistas autorizadas, confirmación, revalidación y escritura.} \\
\bottomrule
\end{tabularx}
\end{table}

\subsection{M5: cinco agentes especializados}

M5 es la arquitectura principal. Sus cinco responsabilidades se coordinan con
un grafo de LangGraph \citep{langgraphGraphApi}, representado en la
Figura~\ref{fig:grafo-m5}:

\begin{figure}[!htbp]
\centering
\begin{tikzpicture}[
  node distance=6mm and 8mm,
  box/.style={draw=ETSBlue,rounded corners,text width=29mm,align=center,
    minimum height=9mm,font=\small,fill=ETSLightBlue!12},
  tool/.style={box,draw=ETSGold!75!black,fill=ETSGold!15},
  flow/.style={-{Latex[length=2mm]},draw=ETSBlue,thick}
]
\node[box] (ctx) {Contexto};
\node[box,right=of ctx] (cmp) {Cumplimiento};
\node[box,right=of cmp] (prf) {Perfiles};
\node[box,below=of prf] (dec) {Orquestación\\(decidir)};
\node[tool,left=of dec] (act) {Actuar\\(herramientas)};
\node[box,below=of dec] (per) {Personalización};
\draw[flow] (ctx) -- (cmp);
\draw[flow] (cmp) -- (prf);
\draw[flow] (prf) -- (dec);
\draw[flow] (dec) to[bend right=15] (act);
\draw[flow] (act) to[bend right=15] (dec);
\draw[flow] (dec) -- (per);
\end{tikzpicture}
\caption{Grafo de M5. Una emergencia o un rechazo terminan antes. El bucle de orquestación está acotado.}
\label{fig:grafo-m5}
\end{figure}

\begin{description}
  \item[Contexto.] Interpreta el último mensaje con la conversación reciente y el
  estado: intención, servicio, fechas, cita o hueco referidos, directivas de
  idioma o detalle, renuncia al perfil y señales clínicas o de urgencia.
  \item[Cumplimiento.] Recomienda continuar, rechazar o escalar y si procede un
  aviso clínico. Su recomendación solo puede restringir: los permisos los
  decide el motor de políticas.
  \item[Perfiles.] Recibe únicamente la vista autorizada para la finalidad y
  decide qué preferencias aplicar a esta petición.
  \item[Orquestación.] Elige la siguiente acción entre las habilitadas por el
  programa para el estado actual. El ciclo de decidir y actuar se repite como
  máximo cuatro veces por turno. Es un límite de seguridad: la gestión más larga
  necesita tres pasos (consultar las citas, buscar huecos y crear la propuesta),
  y el cuarto deja margen. Si se alcanza, el coordinador no sigue llamando a
  herramientas y se responde con lo obtenido hasta ese momento, por ejemplo las
  opciones encontradas, para que la persona continúe en el turno siguiente.
  \item[Personalización.] Redacta el plan de respuesta con el estilo aplicable.
\end{description}

Que haya cinco responsabilidades no implica cinco llamadas por mensaje: el
agente de perfiles solo interviene en reservas, reprogramaciones y
recordatorios, y los mensajes triviales (descartar, rechazar) se redactan sin
modelo.

\subsection{S1: agente único}

S1 integra todas las responsabilidades en un único agente que usa las mismas
herramientas mediante llamadas nativas a funciones. Sus instrucciones recogen
las reglas de confirmación, privacidad, alcance clínico y preferencias. Si su
texto final anuncia un efecto que no ocurrió o pierde una hora relevante, se
sustituye por el texto determinista del último resultado de herramienta.

\subsection{H: flujo programado con tareas LLM}

H fija en código la secuencia de cada gestión: identificar la cita, buscar
huecos, proponer y confirmar. El modelo solo interpreta el mensaje (mismo
agente de contexto que M5) y redacta la respuesta. Sirve de referencia: una
diferencia entre M5 y H no puede atribuirse únicamente al número de agentes,
porque también cambia cuánto del flujo decide el código.

\section{Tratamiento de errores}

\label{sec:lecciones-v1}

En una primera versión del prototipo, los errores de interpretación y los
incumplimientos de seguridad recibían un tratamiento demasiado parecido. Esto
podía terminar una gestión aunque fuera posible aclarar la petición. De esa
experiencia se mantiene una decisión de diseño: comprobar de forma estricta
los permisos y los efectos, y permitir que la conversación resuelva la
información incompleta mediante preguntas o respuestas de respaldo.

\section{Alternativas consideradas}

\paragraph{Grafo explícito frente a bucle propio.} LangGraph representa nodos,
estado y transiciones condicionales con una interfaz común y facilita trazar
qué fase produjo cada decisión. S1 utiliza un bucle de llamadas a herramientas y H sigue una secuencia
programada. Estas dos variantes no necesitan el grafo de M5.

\paragraph{Contratos estructurados frente a texto libre.} Las salidas de los
agentes son objetos JSON validados con Pydantic y enumeraciones cerradas, lo
que impide acciones desconocidas. Los campos no exigen copiar texto del usuario. Los valores dudosos se tratan
como ausentes para que el asistente pueda pedir una aclaración.

\paragraph{Controles comunes frente a controles por variante.} Compartir
herramientas y políticas evita que una variante obtenga permisos distintos y
convierte la seguridad en una propiedad del núcleo, no del modelo. A cambio, la
comparación no mide la aportación aislada de cada barrera determinista.

\paragraph{Persistencia local.} SQLite y una agenda en memoria protegida por
bloqueo permiten reproducir el entorno sin un HIS externo. Las garantías de
concurrencia se limitan a un proceso.
